\chapter{Empirical Runtimes}
\label{appn:runtimes}

All timings were measured on an Intel Core i7-14700HX (20 cores,
single socket) running Python~3.12 with Optuna~3.x in single-threaded
mode.  Wall-clock time was captured with \texttt{time.perf\_counter()}
around each optimiser or simulation call; file I/O and plotting are
excluded.

%% -----------------------------------------------------------------------
\section*{Per-trial cost by optimisation mode}

Table~\ref{tab:rt-per-trial} reports mean per-trial wall-clock time
for each optimisation mode, measured over 20 consecutive trials.

\begin{table}[H]
  \centering
  \caption{Per-trial wall-clock time (20-trial sample, averaged).}
  \label{tab:rt-per-trial}
  \begin{tabular}{@{}lcr@{}}
    \toprule
    Mode & Dimensions & ms\,/\,trial \\
    \midrule
    Transport-only & 5  & 127 \\
    Inventory-only & 25 & 140 \\
    Joint          & 30 & 131 \\
    \bottomrule
  \end{tabular}
\end{table}

Per-trial cost is flat from 5 to 30 dimensions (127--140\,ms).  This
confirms the scalability claim of Section~\ref{sec:opt:dims}: the TPE
bookkeeping overhead is negligible relative to the 365-day simulation
evaluation ($\approx130$\,ms per run), so adding dimensions does not
materially increase cost per trial.

%% -----------------------------------------------------------------------
\section*{Per-experiment runtimes at the canonical trial budget}

The canonical run uses \texttt{--trials 100} for all optimising
experiments.  Table~\ref{tab:rt-per-exp} gives per-experiment
wall-clock times; non-optimising experiments (E0, E1a, E5, E7, E9)
run one or more bare simulation evaluations and are measured directly.
E4 combines an NSGA-II study (\texttt{max(200,\,2$\times$100)=200}
trials) with a constraint sweep (8 targets $\times$ 100 trials = 800
trials), totalling 1000 simulation evaluations.

\begin{table}[H]
  \centering
  \caption{Per-experiment wall-clock time at canonical trial budget (single-threaded,
           \texttt{--trials~100}).}
  \label{tab:rt-per-exp}
  \begin{tabular}{@{}llcrr@{}}
    \toprule
    Experiment & Mode & Dims & Sim.\ evaluations & Wall-clock (s) \\
    \midrule
    E0  & Baseline (1 sim)                  & ---    &    1 &  $<1$ \\
    E1a & Fixed threshold (1 sim)            & ---    &    1 &  $<1$ \\
    E1b & Transport-only                     & 5      &  100 &   13  \\
    E2  & Inventory-only                     & 25     &  100 &   14  \\
    E3  & Joint                              & 30     &  100 &   13  \\
    E3\_per\_sku & Joint (per-SKU)           & 30     &  100 &   13  \\
    E4  & NSGA-II + sweep                    & 30     & 1000 &  130  \\
    E5  & Disruption (2 sims)                & ---    &    2 &  $<1$ \\
    E6  & Policy search (4 $\times$ 100)     & 30--35 &  400 &   53  \\
    E7  & Forecast sweep (5 sims)            & ---    &    5 &    1  \\
    E8  & Bullwhip-aware (5 $\times$ 100)    & 30     &  500 &   65  \\
    E9  & Stochastic (10 sims)               & ---    &   10 &    1  \\
    \midrule
    \textbf{Full suite (E0--E9)} & & & \textbf{$\sim$2220} & \textbf{$\sim$305} \\
    \bottomrule
  \end{tabular}
\end{table}

The full E0--E9 suite completes in approximately \textbf{5 minutes}
on a laptop CPU.  E4 accounts for the largest share ($\sim$130\,s)
because its constraint sweep runs 8 independent joint optimisations in
addition to the NSGA-II study.  E8 ($\sim$65\,s) and E6 ($\sim$53\,s)
follow, each requiring multiple independent optimiser runs.

%% -----------------------------------------------------------------------
\section*{Memory usage}

Peak heap allocation per simulation evaluation, measured with Python's
\texttt{tracemalloc}, is \textbf{11.4\,MB}.  Because Optuna runs
trials sequentially by default, no more than one simulation is live at
a time; total peak process memory (Python interpreter, Optuna state,
and pandas logs) stays below 200\,MB throughout the full suite.
